For \(x\in\mathcal B_{P,\delta}\), write \(S_{P,x}=|D_{P,x}|\) and let \(v_m(u)=(1,u,\ldots,u^{m-1})^\top\).  For arm \(t\) and bandwidth \(0<h\leq r_*\), define
\[
 \widehat G_{t,x}(h)=\frac1{nh^2}\sum_{i=1}^n\mathbf1\{T_i=t\}
 K_\triangle(S_{P,x,i}/h)v_m(S_{P,x,i}/h)v_m(S_{P,x,i}/h)^\top
\]
and
\[
 \widehat r_{t,x}(h)=\frac1{nh^2}\sum_{i=1}^n\mathbf1\{T_i=t\}
 K_\triangle(S_{P,x,i}/h)v_m(S_{P,x,i}/h)Y_i.
\]
Set \(\widehat\mu^{(m)}_{t,h}(x)=e_1^\top\widehat G_{t,x}(h)^{-1}\widehat r_{t,x}(h)\) when the Gram matrix is invertible and set it to zero otherwise; subtract the two arm intercepts.  This is an explicit measurable function of \(U_{n,P}(x)\), so it belongs to \(\mathcal T_n^{\mathrm{dist}}\).

The coarea identity and the assumed radial mass bounds give
\[
 \mathbb E_P\widehat G_{t,x}(h)
 =\int_0^1K_\triangle(u)v_m(u)v_m(u)^\top\frac{Q_{t,P,x}(hu)}h\,du.
\]
Consequently, for every \(a\in\mathbb R^m\),
\[
 \underline c\int_0^1uK_\triangle(u)\{a^\top v_m(u)\}^2du
 \leq a^\top\mathbb E_P\widehat G_{t,x}(h)a
 \leq\overline c\int_0^1uK_\triangle(u)\{a^\top v_m(u)\}^2du.
\]
The middle moment matrix is positive definite because a nonzero polynomial cannot vanish on an interval of positive Lebesgue measure.  Since \(m\) is fixed, its smallest eigenvalue is a strictly positive constant.  This proves all population rank, positivity, and inverse conditions.

By the stated coarea remainder, the radial regression equals a polynomial of degree at most \(m-1\) whose intercept is \(\mu_{t,P}(x)\), plus a remainder bounded by \(As^m\).  The local-polynomial normal equations reproduce the polynomial exactly, for every empirical design on which the Gram matrix is invertible.  The population inverse bound therefore gives intercept bias at most \(Ch^m\).

The smooth-boundary length lemma supplies a polynomial-size \(n^{-3}\)-net of \(\mathcal B_{P,\delta}\).  Each Gram summand is bounded after scaling and is nonzero with probability at most \(Ch^2\), by \(Q_{t,P,x}(s)\leq\overline c s\).  Matrix Bernstein on the net and the Lipschitz dependence of \(K_\triangle(S_{P,x}/h)v_m(S_{P,x}/h)\) on \(x\) imply
\[
 \sup_{P,t,x}\left\|\widehat G_{t,x}(h)-\mathbb E_P\widehat G_{t,x}(h)\right\|
 =O_P\!\left(\sqrt{\frac{\log n}{nh^2}}\right).
\]
At the stated bandwidth, \(nh^2/\log n\to\infty\).  Weyl's inequality therefore makes every sample Gram matrix invertible with uniformly bounded inverse outside an event of probability at most \(n^{-4}\).

Conditionally on the signed radius, \(Y-g^{\mathrm{dist}}_{t,P,x}(S_{P,x})\) is uniformly sub-Gaussian: it is the sum of the conditional sub-Gaussian regression error and a centered angular regression fluctuation bounded by the H\"older envelope.  The corresponding conditional exponential inequality, the same net, and Lipschitz extension yield
\[
 \sup_{P\in\mathcal C}\mathbb E_P\sup_{x\in\mathcal B_{P,\delta}}
 |\widehat\tau^{(m)}_{n,h}(x)-\tau_P(x)|
 \leq C\left\{h^m+\sqrt{\frac{\log n}{nh^2}}\right\}.
\]
On singular-matrix events the estimator is zero; the target is uniformly bounded, so their expected contribution is negligible.  Balancing the two terms gives
\[
 h\asymp\left(\frac{\log n}{n}\right)^{1/(2m+2)},
 \qquad
 R_{n,\delta}^{\mathrm{dist}}(\mathcal C)
 \leq C\left(\frac{\log n}{n}\right)^{m/(2m+2)}.
\]

For the conditional reverse inequality, use the uniform-design straight-interface H\"older bump submodel of order \(m\) whose containment in \(\mathcal C\) is the theorem's explicit premise.  Place disjoint bumps of width \(h\) and height \(ch^m\) at \(M\asymp h^{-1}\) trimmed centers and use bounded Gaussian-noised Bernoulli outcomes.  The premise supplies submodel membership; directly, adjacent target values differ by \(ch^m\), and the score mass of one bivariate bump is \(O(h^2)\), so
\[
 \operatorname{KL}(P_j^{\otimes n},P_0^{\otimes n})\leq Cnh^{2m+2}.
\]
At the displayed bandwidth, choosing the fixed amplitude small makes this at most \(\alpha\log M\) for some \(\alpha<1\).  For a uniform bump index \(J\),
\[
 I(J;O_{1:n})\leq M^{-1}\sum_j\operatorname{KL}(P_j^{\otimes n},P_0^{\otimes n}),
 \qquad
 I(J;O_{1:n})\geq(1-p_e)\log M-\log2.
\]
Thus every decoder has error probability bounded away from zero.  Sup-norm accuracy below half the bump height would decode the center, proving the matching lower bound \(ch^m\).  Exactness is therefore conditional on the stated submodel containment, as claimed.

If a class instead contains the shrinking angular-imbalance family generated by \(\mathfrak H_{\mathrm{ang}}\), its adjacent target separation is \(ca_n\).  In the mean-\(2n\) Poisson experiment, the disjoint local processes are conditionally independent standard-Borel channels and the outside process is bit-independent.  The Kullback--Leibler slack from the angular packing and the coordinatewise-overlap direct-product lemma give probability at least \(1/2-o(1)\) of some bit error, even after all cross-center observations are revealed.  Thresholding and exponentially accurate subsampling de-Poissonization yield the lower frontier \(ca_n\).  For \(m=1\), the upper rate equals \(a_n\), in agreement with smooth-geometry saturation; for \(m\geq2\), \(m/(2m+2)>1/4\).

Finally, the two contrasting constructions establish the operational claim.  Bounded continuity of \(f_P\) alone permits shrinking angular imbalance and hence does not suffice for a faster order.  Conversely, density differentiability is not necessary for the displayed product-integral remainder: if the arm regression is locally constant on each relevant circle, then \(N_{t,P,x}(s)/Q_{t,P,x}(s)=\mu_{t,P}(x)\) for any bounded continuous positive density, irrespective of density derivatives.  Thus the quantities actually used by the estimator and its Gram matrix are the radial mass bounds and the normalized product-integral remainder, not regularity of the density in isolation.  This conclusion concerns the identifying functional while keeping it distinct from the causal target \(\tau_P\).